\section{Dynamics-Encoding Bi-Implication}
\label{app:hom-biimp}

This appendix proves the bi-implication stated in Section~\ref{sec:dynamics-encoding-fragment}.

\paragraph{Statement.}
$Z_{\text{impl}}$ satisfies the hom-relative dynamics-encoding property set compiled from $Z_{\text{spec}}$
if and only if
there exists a Wymore system homomorphism $h : Z_{\text{impl}} \to Z_{\text{spec}}$ (Def.~4.3).

\subsection{Setup}
\label{app:hom-setup}
Write $Z_{\text{spec}}=(S_{\text{spec}},I_{\text{spec}},O_{\text{spec}},\delta_{\text{spec}},\rho_{\text{spec}})$ and $Z_{\text{impl}}=(S_{\text{impl}},I_{\text{impl}},O_{\text{impl}},\delta_{\text{impl}},\rho_{\text{impl}})$. Types may differ. A run of $Z_{\text{impl}}$ is determined by $s_0$ and $f:T\to I_{\text{impl}}\cup\{\mathrm{none}\}$ via \eqref{eq:trace-init}--\eqref{eq:trace-readout}. Maps $h_S,h_I,h_O$ are surjective; $h_I$ lifts to options by sending $\mathrm{none}$ to $\mathrm{none}$.

The admissible class quantifies over all states of $Z_{\text{impl}}$ as initial states and all input trajectories. This is used in one direction only---recovering the static laws from the runs---and it is necessary there: a state reached by no admissible run is unconstrained by any run-based reading, whereas Definition~4.3 constrains every state. When a designated initial state is intended, replace $Z_{\text{impl}}$ throughout by the subsystem it generates; the argument is unchanged and the conclusion is about that subsystem.

\subsection{Hom-relative fragment}
\label{app:hom-fragment}
Clauses are indexed by specification states. The atom $\mathit{specState}(s)$ holds at tick $t$ when $h_S(g(t))=s$. For each spec state $s$:
\begin{enumerate}[label=(\roman*), leftmargin=2em]
    \item Open readout: if $\rho_{\text{spec}}(s)=o$ and $\mathit{specState}(s)$, then $h_O(y(t))=o$.
    \item Closed readout: if $\rho_{\text{spec}}(s)=\mathrm{none}$ and $\mathit{specState}(s)$, then $h_O(y(t))=\mathrm{none}$.
    \item Autonomous: if $\mathit{specState}(s)$ and $f(t)=\mathrm{none}$, then $h_S(g(t{+}1))=\delta_{\text{spec}}(s,\mathrm{none})$.
    \item Guided: if $\mathit{specState}(s)$, $f(t)=i_{\text{impl}}$, and $h_I(i_{\text{impl}})=i_{\text{spec}}$, then $h_S(g(t{+}1))=\delta_{\text{spec}}(s,i_{\text{spec}})$.
\end{enumerate}
Write $Z_{\text{impl}}\models_{\mathrm{dyn}}\Phi^{\mathrm{hom}}_{Z_{\text{spec}}}$ when there exist surjective $h_S,h_I,h_O$ such that every admissible run satisfies every clause.

\subsection{Def.~4.3 homomorphism}
\label{app:hom-def}
A homomorphism requires, for every impl state $x$,
\begin{align}
  h_O(\rho_{\text{impl}}(x)) &= \rho_{\text{spec}}(h_S(x)) \label{eq:hom-readout-static}\\
  h_S(\delta_{\text{impl}}(x, oi)) &= \delta_{\text{spec}}(h_S(x), oi.\mathit{map}(h_I)). \label{eq:hom-step-static}
\end{align}

\subsection{Bridge lemma}
\label{app:hom-bridge}
Fix surjective $h_S,h_I,h_O$. The static laws~\eqref{eq:hom-readout-static}--\eqref{eq:hom-step-static} hold at every impl state if and only if every impl run satisfies the four clause families.

\paragraph{Static $\Rightarrow$ runs.} Trajectory equations give $g(t{+}1)=\delta_{\text{impl}}(g(t),f(t))$ and $y(t)=\rho_{\text{impl}}(g(t))$. Instantiating the static laws at $g(t)$ yields each clause.

\paragraph{Runs $\Rightarrow$ static.} From impl state $x$, run with constant input $oi$ (or $\mathrm{none}$). The clause at $t=0$ recovers the static law at $x$.

\subsection{Directions}
\label{app:hom-forward}
\label{app:hom-backward}
If $h$ is a homomorphism, the static laws hold, hence by the bridge lemma every run satisfies the clauses. Conversely, if $Z_{\text{impl}}\models_{\mathrm{dyn}}\Phi^{\mathrm{hom}}_{Z_{\text{spec}}}$, the witnessing maps satisfy the static laws at every state and are therefore a Def.~4.3 homomorphism. Surjectivity ensures every specification row is inhabited by some impl state.

\subsection{Theorem and fixed-table corollary}
\label{app:hom-fixed}
\[
  Z_{\text{impl}} \models_{\mathrm{dyn}} \Phi^{\mathrm{hom}}_{Z_{\text{spec}}}
  \;\;\Longleftrightarrow\;\;
  \exists\, h : Z_{\text{impl}} \to Z_{\text{spec}}\ \text{Wymore homomorphism}.
\]
When types coincide and identity maps are used, this specializes to pointwise $\delta$/$\rho$ agreement (partial identity homomorphic image). Formulas outside the fragment---including $F$ properties and output-table-only encodings---are outside the bi-implication. Global FC equality for arbitrary such $\Phi$ remains open.

\paragraph{Machine-checked formalization.}
This bi-implication, the invariance and composition results of Section~\ref{sec:invariance}, and every conformance and rejection claim of the case study are formalized in a proof assistant, including the coupling recipe of Section~\ref{sec:case-machine} and its resultant, the accepted and rejected zone variants, the lift of zone conformance to the coupled system, and the tick-granularity negative result. The rejections are proofs that no maps exist rather than reports of failed search. The finite instances decided by the checker of Section~\ref{sec:case-solver} are compared against those theorems wherever a theorem covers the instance, and one map found by the checker is re-verified inside the formalization.
